\section{A canonical mass on the trace}
\label{mass:section}

We now turn the Green potential into an increasing measure on the trace.
The conditional Green-potential approach goes back to
Lawler--Sheffield \cite[Proposition~2.2 and Section~3]{LawlerSheffield2011}
and Lawler--Zhou \cite[Proposition~A and Theorem~1]{LawlerZhou2013};
see also the Minkowski-content construction in \cite{LawlerRezaei2015}.
Here one common sequence of convex combinations supplies the mass and
all of its versions. We construct the measure and prove the
properties needed below directly; no identification with Minkowski content
is required.  All measures in this section are measures on the open domain.
Write $\lambda(f)=\int f\,\dd\lambda$ for integration against a measure.

\begin{proposition}[Construction and conditional kernel]
\label{mass:construction}
There is an adapted increasing family $(\mu_t)_{0\le t\le\infty}$ of
locally finite measures on $\HH$, starting at zero, with the following
properties.  Its terminal measure is nonatomic and has moments of every
integer order on compact sets.  Almost surely, simultaneously for all
$0\le v<t<\infty$, the increment $\mu_t-\mu_v$ is carried by
$\gamma([v,t])\cap\HH$.  The process is locally continuous in total
variation, and $\mu_t\uparrow\mu_\infty$ as $t\to\infty$.
For every stopping time $S$, possibly infinite,
\begin{equation}
 \EE\bigl[\mu_\infty(\dd x)-\mu_S(\dd x)\mid\mathcal F_S\bigr]
   =G_S(x)\,\dd A(x),
 \qquad G_\infty:=0.
 \label{mass:kernel}
\end{equation}
This is an equality of conditional measure kernels; in particular it
allows nonnegative tests known at time $S$.  Moreover,
\begin{equation}
 \EE\mu_\infty(\dd x)=G_0(x)\,\dd A(x).
 \label{mass:mean}
\end{equation}
\end{proposition}

\subsection{Deterministic convex extraction}

The compact moment bounds permit one simultaneous limit for all
rational times and a determining family of spatial tests. Keeping the
same convex coefficients is what later makes the kernels and conformal
images properties of one measure rather than separately chosen versions.

For an integer $j\ge1$ define
\begin{equation}
 M_t^{(j)}(\dd x)
 =G_{\tau_j(x)}(x)
   \ind_{\{\tau_j(x)\le t,\ \tau_j(x)<\infty\}}\,\dd A(x),
 \qquad 0\le t\le\infty.
 \label{mass:approximants}
\end{equation}
The density is zero on failure to reach the level.  For a compact
$C\Subset\HH$, fix $c_C>0$ with $c_C\le\inf_C R_0$.  Given distinct
$x_1,\ldots,x_m\in C$, set
\[
 b_i=\min\bigl(c_C,\min_{k\ne i}|x_i-x_k|\bigr),
\]
where the second minimum is omitted when $m=1$.
The terminal approximant is the level weight $A_j$ of
Corollary~\ref{mp:level-moments}.  The collision-kernel estimate in its
proof gives, for every bounded nonnegative $f$ supported in $C$,
\begin{equation}
 \sup_j\EE\bigl[M_\infty^{(j)}(f)^m\bigr]
 \le C_{m,a}\int_{C^m}
       \prod_{i=1}^m f(x_i)b_i^{-\alpha}
       \,\dd A(x_1)\cdots\dd A(x_m)<\infty.
 \label{mass:moment-bound}
\end{equation}
The last assertion is Lemma~\ref{mp:kernel-integral}.
All constants are independent of $j$.

We obtain one sequence of convex combinations as follows.  Choose a
countable family of continuous compactly supported tests, uniformly
dense among tests supported in each member of a compact exhaustion.
Include nonnegative cutoff functions equal to one on that exhaustion,
and rational linear combinations of the chosen tests.  Enumerate the
pairs of these tests with times in $\mathbb Q_+\cup\{\infty\}$ as
$(f_r,t_r)$, $r\ge1$, and set
$X_r^{(j)}=M_{t_r}^{(j)}(f_r)$.
Each coordinate is bounded in $L^2$ by \eqref{mass:moment-bound}.
A diagonal subsequence $j_l$ has a weak limit $X_r$ in every coordinate.

For each $k$, consider the first $k$ coordinates as a vector in
$(L^2)^k$.  Its weak limit $(X_1,\ldots,X_k)$ lies in the norm closure
of the convex hull of every tail.  Hence there are deterministic
coefficients $c_{kl}\ge0$, supported on finitely many $l\ge k$, with
$\sum_lc_{kl}=1$, such that
\[
 \sum_{r=1}^k\left\|
       \sum_{l\ge k}c_{kl}X_r^{(j_l)}-X_r
                    \right\|_2^2\le4^{-k}.
\]
Define
\[
 \overline M_t^{(k)}=\sum_{l\ge k}c_{kl}M_t^{(j_l)}.
\]
For every listed pair, the summable squared errors give both $L^2$
and almost-sure convergence; intersecting the resulting events over
$r$ gives simultaneous convergence on the whole countable list.
The same coefficient row is used for every time and spatial test.

On this event, the cutoff tests bound all local masses.  Density therefore
extends convergence to every continuous compactly supported test.
The limiting positive functionals are Radon measures, with almost-sure
vague convergence at all listed times.  Positivity and ordering in time
pass to the limit.  The limits at finite rational times are adapted, since
the corresponding $L^2(\mathcal F_t)$ subspaces are closed.  Convexity and
Fatou's Lemma preserve \eqref{mass:moment-bound} for the terminal measure
and continuous nonnegative tests.  Approximation then gives the compact
moment assertions.

We also obtain nonatomicity, which will be needed to pass from rational
times to all times.  For two points in a fixed larger compact the kernel
in \eqref{mass:moment-bound} is bounded by
$C(1+|x-y|^{-2\alpha})$.  Cover a given smaller compact by mesh squares
of side $\delta$, and majorize each square's indicator by a continuous
cutoff in a slightly enlarged square.  The enlarged squares have bounded
overlap.  The expected sum of their squared masses is bounded by
\[
 C\iint_{|x-y|\le C\delta}
       (1+|x-y|^{-2\alpha})\,\dd A(x)\dd A(y),
\]
where both variables lie in the larger compact.  This tends to zero
because $2\alpha<2$.  Every atom in the smaller compact contributes the
square of its mass to the corresponding sum at every mesh size.
Fatou's Lemma, followed by a countable compact exhaustion, proves that
$\mu_\infty$ has no atoms.

\subsection{Support, continuity, and stopping times}

The finite-level conditional weight identity of
Lemma~\ref{op:finite-tilt} gives, at a finite rational time $t$,
\begin{equation}
 \EE\bigl[M_\infty^{(j)}(\dd x)-M_t^{(j)}(\dd x)
             \mid\mathcal F_t\bigr]
 =G_t(x)\ind_{\{x\in H_t,\ R_t(x)>R_0(x)e^{-2aj}\}}\,\dd A(x).
 \label{mass:approx-kernel}
\end{equation}
The strict inequality excludes a level already attained at $t$.
At every surviving point the indicator tends to one.  Since
$\EE G_t(x)\le G_0(x)$, dominated convergence on compact tests applies to
the right side.  Local $L^2$ bounds give uniform integrability on the
left.  Passing through the tail convex combinations proves
\eqref{mass:kernel} at rational times.  The mean-one finite-level weights
also give \eqref{mass:mean}.

These identities initially hold for countably many deterministic tests.
They identify the conditional measure kernel, and hence, by the Monotone
Class Theorem, hold against every nonnegative
$\mathcal F_t\otimes\mathcal B(\HH)$-measurable test.  In particular,
testing the complement of the known surviving domain shows that
\begin{equation}
 \mu_\infty-\mu_t\text{ is carried by }H_t
 \quad\text{at every finite rational }t,
 \label{mass:future-support}
\end{equation}
on one event of probability one.

At the same times, $\mu_t$ is carried by the past trace.  Indeed, on a
fixed compact, a point receiving mass in $M_t^{(j)}$ has, at its level
time, distance at most $R_0(x)e^{-2aj}$ from the current domain boundary.
For large $j$ that nearby boundary cannot be the real line, so it belongs
to $\gamma([0,t])$.  Thus the approximants, and their tail convex
combinations, are carried locally by shrinking neighborhoods of this
closed trace segment.  Vague convergence gives the claimed support.
For rational $v<t$, the increment is dominated by $\mu_t$ and by
$\mu_\infty-\mu_v$, so it is carried by $\gamma([v,t])\cap\HH$.

Take increasing rational left limits and decreasing rational right limits
of the measure process.  At any finite time $w$, their difference is
dominated by increments over every rational interval containing $w$.
Continuity of the trace forces its support to lie in the singleton
$\{\gamma(w)\}\cap\HH$.  The difference vanishes by nonatomicity.
At zero use zero as the left limit; the same argument gives zero, since
$\gamma(0)\notin\HH$.  This argument is pathwise on the one event already
constructed and therefore works simultaneously for every $w$.
It defines an increasing process on all times, locally continuous in total
variation: on a compact the variation of a positive increment is just its
mass, and the possible one-sided jump measures have been shown to vanish.
Adaptedness follows from the rational left limits.
The support assertion for arbitrary intervals follows by enclosing them
in rational intervals and taking limits.

The terminal value is exhausted by finite times.  By
Proposition~\ref{op:potential}, $\EE G_t(x)\to0$ as $t\to\infty$.
Apply \eqref{mass:kernel} at integer times and use
$\EE G_t\le G_0$ on compact tests.  The nonnegative measure
$\mu_\infty-\lim_{t\to\infty}\mu_t$ has zero mean, and hence is zero.
This also gives local total-variation convergence at infinity.

For completeness, let $S$ be any stopping time.  On the event of rational
convergence, monotonicity squeezes $\overline M_S^{(k)}(f)$ between its
values at rational times approaching each realized finite $S$.
Continuity gives convergence to $\mu_S(f)$ for every nonnegative
continuous compact test.  On $\{S=\infty\}$ use terminal convergence of
the extracted sequence.  Uniform integrability again follows from the
terminal $L^2$ bound.  The finite-level identity
\eqref{mass:approx-kernel} holds with $t$ replaced by $S$, and with zero
right side on $\{S=\infty\}$.  The bound
$\EE[G_S(x);S<\infty]\le G_0(x)$ from
Proposition~\ref{op:potential} permits the same dominated convergence.
This proves \eqref{mass:kernel} at $S$, including its conditional-kernel
interpretation, and completes Proposition~\ref{mass:construction}.

\begin{corollary}[First-visit accounting]
\label{mass:first-visit}
Put $v(x)=\inf\{t\ge0:\gamma(t)=x\}$, with value infinity if the set is
empty.  Almost surely $v(x)<\infty$ for $\mu_\infty$-almost every $x$, and
simultaneously for every finite $t$,
\begin{equation}
 \mu_t(\dd x)=\ind_{\{v(x)\le t\}}\mu_\infty(\dd x).
 \label{mass:first-visit-identity}
\end{equation}
\end{corollary}

\begin{proof}
At rational $t$, past mass is carried by $\gamma([0,t])$, while remaining
mass is carried by the open set $H_t$, disjoint from that past trace.
This proves the identity there.  Increasing finite times exhaust the
terminal measure.  Right approximation by rational times, continuity of
the measure process and of the trace, then proves the identity at all
times.  In particular a later geometric return to an already visited
point introduces no new mass there.
\end{proof}

\subsection{Versions, conformal maps, and restart}

We have constructed the continuous mass and its conditional kernel.
To use them at centered stops and in random future domains, we next
fix a measurable path functional and establish equality of the actual
restarted measures.

We fix one nonanticipating measurable version of this construction.
On the canonical continuous-path space each $M^{(j)}$, tested on an
interior compact, is a raw adapted c\`adl\`ag process: its terminal density
has a deterministic finite bound there, and one-sided continuity follows
from dominated convergence of the hitting-time indicators.
The same holds for every deterministic convex combination.
At each finite time, take the vague limit of $\overline M_t^{(k)}$
where it exists, and use the zero measure otherwise.  The space of
locally finite positive Radon measures is Polish in the vague topology;
its convergence set and limit map are therefore Borel.  This defines
a raw progressively measurable measure-valued functional.  At infinity use the same limit rule for
$\overline M_\infty^{(k)}$.  Rational squeezing as above shows that it
agrees, simultaneously at every time under the ordinary Brownian law,
with the continuous process just constructed.

All evaluations under a centered law use this raw functional of the
stopped path.  At each finite inner level the density in
Lemma~\ref{op:finite-tilt} transfers the required statements about
that stopped piece.  Consistent laws of the stopped pieces, and increasing
limits of their coordinates, suffice thereafter.  We do not extend the
centered law to the original Brownian space completed by infinite-future
null sets.  This distinction also ensures parameter measurability:
Loewner solutions, level times and conformal maps are measurable in the
driver and marked point, and progressive evaluation of the fixed mass
functional at those times is measurable.

For a fixed conformal identification $J:\HH\to D$, define
\begin{equation}
 \mu_t^D=J_*\bigl(|J'|^d\mu_t\bigr).
 \label{mass:conformal}
\end{equation}
This rule holds exactly for the approximants: conformal radius contributes
$|J'|^{-\alpha}$ and area contributes $|J'|^2$, so their product is
$|J'|^d$.  Inner-clock ratios are unchanged.
A compact continuous test in $D$ pulls back to such a test in $\HH$.
Thus the same deterministic convex combinations converge in every fixed
image domain, on the event giving convergence for all compact continuous
tests.  The kernel, local continuity, nonatomicity and support statements
hold there with the corresponding $G_t^D$. More precisely, its increments
are carried by $J(\gamma([v,t])\cap\HH)$. The map is applied only to
interior trace points; an arbitrary $J$ need not extend continuously
to the boundary.
The construction is jointly measurable in a measurable family of maps
$J$; it requires no new extraction for different maps.

The kernel also proves uniqueness.  If $A_t$ and $\widetilde A_t$ are
two continuous increasing adapted test-mass processes starting at zero,
with integrable terminal values and the same remaining conditional
potential, then
\[
 A_t-\widetilde A_t
   =\EE[A_\infty-\widetilde A_\infty\mid\mathcal F_t].
\]
Their difference is a continuous martingale of finite variation, hence
vanishes.  Applying this to a countable determining set of tests proves
uniqueness of the measure process.  In particular it removes dependence
on the convex extraction.  It also gives Brownian scaling: if
$B^{(c)}_t=c^{-1}B_{c^2t}$, then, for each $c>0$,
\begin{equation}
 \mu_t(B^{(c)})(A)=c^{-d}\mu_{c^2t}(B)(cA)
 \quad\text{almost surely, for all }t\text{ and Borel }A\subset\HH.
 \label{mass:scaling}
\end{equation}
Indeed the scaled process has the same kernel, continuity and terminal
integrability as the canonical process for $B^{(c)}$.

For natural parametrization, Lawler--Rezaei
\cite[Proposition~3.13]{LawlerRezaei2012Basic} prove conformal restart
at an almost surely finite stopping time, simultaneously over later
time intervals.  We prove the following measure-valued identity for
the mass constructed here directly from its conditional kernel.

\begin{proposition}[Pathwise interior restart]
\label{mass:restart}
Fix a stopping time $S$.  On $\{S<\infty\}$ put
$J_S=(g_S-B_S)^{-1}:\HH\to H_S$, and let $\widetilde\mu$ be the canonical
mass functional of the fresh Brownian motion $B_{S+t}-B_S$.
Almost surely, as measures on $H_S$, simultaneously for all
$0\le t\le\infty$,
\begin{equation}
 (\mu_{S+t}-\mu_S)|_{H_S}
       =(J_S)_*\bigl(|J_S'|^d\widetilde\mu_t\bigr).
 \label{mass:restart-identity}
\end{equation}
The same assertion holds when the initial domain is a fixed conformal
image of $\HH$.  Once this equality of measures has been obtained,
one may integrate any random nonnegative Borel weight on the restart
domain and evaluate at any later random time.
\end{proposition}

\begin{proof}
Write $\nu_t$ for the right side of \eqref{mass:restart-identity} on
$\{S<\infty\}$, and set both processes in the comparison below to zero
on $\{S=\infty\}$.  Take a countable collection of rational balls
$U\Subset\HH$, and dense
countable families of nonnegative tests $\phi\in C_c(U)$, including local
cutoffs.  Localize each comparison to
\[
 E_U=\{S<\infty,\ \overline U\subset H_S\}\in\mathcal F_S.
\]
We first check terminal integrability, before invoking uniqueness.
Conditioning on $\mathcal F_S$, the strong Markov property, the canonical
mean measure and change of variables give
\begin{equation}
 \EE[\ind_{E_U}\nu_\infty(\phi)\mid\mathcal F_S]
       =\ind_{E_U}\int\phi(x)G_S(x)\,\dd A(x).
 \label{mass:restart-integrability}
\end{equation}
The pulled-back test is random, with potentially random derivative and
support bounds.  The formula remains valid: first use nonnegative simple
functions of the restart data and source point, and then apply conditional
Tonelli and monotone convergence.  Its expectation is at most
$\int\phi G_0\,\dd A<\infty$.  Thus the candidate terminal test mass is
integrable.  The actual terminal increment is bounded by
$\mu_\infty(\phi)$ and is integrable as well.

At $S+t$, the candidate and actual increment have the same remaining
conditional potential.  For the candidate, this follows from the
canonical kernel for the fresh driver, with its $\mathcal F_S$-measurable
pulled-back test; the Loewner composition rule and
\eqref{mass:conformal} transform that potential into $G_{S+t}$.
For the actual process it is \eqref{mass:kernel} at $S+t$.
Consequently
\[
 D_t:=\ind_{E_U}\bigl[(\mu_{S+t}-\mu_S)(\phi)-\nu_t(\phi)\bigr]
       =\EE[D_\infty\mid\mathcal F_{S+t}].
\]
This continuous finite-variation martingale starts at zero and is
identically zero.  Intersect the resulting events over the countable
balls and tests, and use continuity to obtain equality for all times.
Those balls whose closures lie in the realized $H_S$ cover it; their
tests determine its locally finite measures.  This proves the pathwise
measure equality throughout the random domain.

Only now do we integrate random weights.  In particular, if
$F:\DD\to H_S$ is a random conformal map and $K\Subset\DD$, the Borel
weight
\[
 \ind_{F(K)}(x)\,|F'(F^{-1}(x))|^{-d}
\]
is bounded on its realized compact support.  Its integration follows
directly from equality of the realized measures, without a further
exceptional event indexed by $F$ or $K$.  The all-time equality similarly
permits substitution of a later random time.  Composition with a fixed
initial identification proves the last domain assertion.
\end{proof}

To obtain the centered restart rule, first work under the ordinary law.
For a stopping time $S$, work on $\{S<\infty,\ z\in H_S\}$ and write
$g_S(z)-B_S=\lambda e^{i\theta}$, where $\lambda>0$.
Conditionally on $\mathcal F_S$, the rescaled fresh driver
$\lambda^{-1}(B_{S+\lambda^2u}-B_S)$ is standard Brownian motion.
We may apply \eqref{mass:scaling} at this random scale: after conditioning,
$\lambda$ is fixed, the fresh driver is independent of $\mathcal F_S$,
and the scaling equality holds almost surely at each fixed scale.
Joint measurability of the mass functional permits this conditioning.
Together with Proposition~\ref{mass:restart}, scaling identifies the
actual future mass in these coordinates.

Now apply the finite change of law through the next radius target of
$z$.  The ratio of Green weights is conformally invariant, since the
derivative at the marked point cancels between the two weights.
Lemma~\ref{op:finite-tilt} therefore identifies, conditionally on the
stopped past under the centered law, the future mass pulled back to a
disk centered at $z$ with the canonical disk law of starting angle
$\theta$.  This assertion uses only data through the finite test endpoint.

\begin{proposition}[Positive mass on every interval]
\label{mass:positive}
Almost surely, for every \mbox{$0\le v<t<\infty$},
\[
 (\mu_t-\mu_v)(\HH)>0.
\]
The positive mass can always be detected on an interior compact.  The
same statement holds in every fixed conformal image, in its associated
half-plane time.
\end{proposition}

\begin{proof}
Equation~\eqref{mass:mean} gives positive terminal expectation on a
compact with nonempty interior.  Terminal exhaustion therefore gives
positive probability of positive interior mass by some finite
deterministic time.  Scaling shows that
$\PP\{\mu_t(\HH)>0\}$ has the same positive value for every $t>0$.
The raw event $\limsup_{n\to\infty}\{\mu_{1/n}(\HH)>0\}$ belongs
to the Brownian germ sigma field. On the common event of monotonicity it
equals $\bigcap_{n\ge1}\{\mu_{1/n}(\HH)>0\}$, so has this same
probability. The germ zero--one law
makes it one.  Hence positive mass is present before every positive time.
Restart at each rational time and apply
Proposition~\ref{mass:restart}.  The strictly positive conformal weight
preserves positivity, giving positive increments on all rational
intervals on one event.  Every nontrivial real interval contains a
nontrivial rational interval.  A countable compact exhaustion detects
the positive increment mass; \eqref{mass:conformal} proves the image-domain
assertion.
\end{proof}
